% !TeX root = ../papers/nc-p3.tex
\section{純テンソル関係のない曲線射影での全分岐の排除}
\label{sec:final-pure-branch-frames}
\begin{theorem}\label{thm:final-pure-branch-frames}
標準正則窓で全ての射影像が高々一次元なら、片側分岐があるためには少なくとも一辺に
純テンソル関係が必要である。従って純関係自由な曲線射影の全窓では、反対階数によらず、
全辺の両側のファイバーは空か一点であり、点グラフの場合の判定が適用できる。
\end{theorem}
\begin{proof}[共通の基底]
前定理の二配置だけを扱う。位置 $j$ において初回の直線を $L_j$、終回の直線を $M_j$ と書く。
これらがともにあるときは一点 $e_j$ で交わる。
実際、それらを与える隣の順・逆中心の縮約は非零である（順中心の中央支持が4）。
その位置で両直線の消滅空間の交叉に支持されるので直線は交わり、
一致すれば間に一つだけ動く位置を持つ列となってしまう。
$L_j$ の全点は前点が固定した一つなので逆中心を含まず、逆に $M_j$ は順中心を含まない。
二直線列の初回・終回にある順・逆中心は系\ref{cor:two-line-secant-roots}で与えられる。

四中心の場合、各位置の二直線を $L_j(a,b)=(a,b,0,0), M_j(a,b)=(0,a,b,0)$、
順・逆中心を座標点 $0,2$ とすることができる。
ここで座標点の番号は $0,1,2,3$ とし、$e_j$ は1である。
実際、移動 $L_j\to M_{j+1}$ は $F_j$ を $e_{j+1}$ に、$e_j$ を $B_{j+1}$ に送るから、
各移動でそれ以外のパラメータも揃えられる。交わる二直線の基底は交点の代表の倍率を合わせて同時に延ばせる。
三中心の場合には位置0の三次曲線を $C(a,b)=(a^3,a^2b,ab^2,b^3)$ とし、
そこでの順中心を $C(1,0)$、逆中心を $C(0,1)$ とする。両中心がこの曲線にあることは列の次数公式による。
位置 $1,2,3$ の二直線は上と同じ形で、$L_3\to C\to M_1$ のパラメータを揃える。
$e_3$ は逆中心 $C(0,1)$ に移り、$e_1$ は順中心 $C(1,0)$ から来るので、
残り二移動 $L_1\to M_2, L_2\to M_3$ と整合する。

各文字で点座標の最初の三座標だけを使う制限を三座標制限と呼ぶ。
$L_j\to M_{j+1}$ の辺では、$F_j\times L_{j+1}, M_j\times B_{j+1}$ を合わせると
三座標制限後の関係は全て次の形に入る。
\begin{equation}\label{eq:pure-branch-core-four}
         \begin{pmatrix}0&0&c\\ d&-c&0\\ h&f&0\end{pmatrix}.
\end{equation}
また、二直線列の最初の中心 $F_j$ に対して $R_j$ の第二点を $L_{j+1}$ に制限した射影は
$\ker F_j\otimes H^0(L_{j+1},\mathcal O(1))$ への同型である。
逆に終回の直線による反対向きの制限も同型である（命題\ref{prop:two-line-branch-block}の束の証明）。
\end{proof}

\begin{proof}[四中心の排除]
全位置を三座標に制限したテンソルを $W$ とする。その係数は各隣接対で式\eqref{eq:pure-branch-core-four}
を満たし、解空間は一次元である。$W\ne0$ である。実際、一つの二直線列上で窓は括弧と階数3の完全対合の積であり、
その残り二文字を三座標に制限する際に対合の両側で各一次元しか失わず、階数は少なくとも1である。
倍率を定め、行・列の順序を対 $10,11,20,21$ とすると、各切断での四次小行列は
\[
 K=\begin{pmatrix}0&0&0&1\\0&-1&0&0\\0&0&-1&0\\1&0&0&0\end{pmatrix}.
\]
一位置だけ第四座標3、他を三座標とした許容係数の空間は次元4である。位置0をその位置とすると
$W$ の位置0に三座標のベクトルを代入した三行と、単一語 $3210$ の係数だけでこれを張る。
これらの小さい線型計算は各対 $00,01,12,22$ の係数を零、$02$ と $11$ の係数の和を零として得る。
添付コードは有理数上の核と基底を確認する。
従って各位置で第四の点座標ベクトルを三座標分だけ平行移動し、
一位置だけ3の係数を $3210$ の巡回語（位置ごとに独立な倍率）のみにできる。これらの変更は同時に行える。

位置0,2を3に固定した三行三列を $N_{ab}=w_{3a3b}$ とする。$(0,1)\mid(2,3)$ 切断の
$K$ の四行四列を消去してもこの部分は $N$ のままである。
$3a$ 行では $a=2$ のときだけ $10$ 列が残り、$3b$ 列でも $10$ 行しか使わず、
$(K^{-1})_{10,10}=0$ だからである。他の $(1,2)\mid(3,0)$ 切断でも $a3$ 行、$b3$ 列で
修正に使う対はともに $21$ で、再び逆行列の対角係数が零である。
第一の切断における $N$ の $a,b\in\{0,1\}$ 部分は可逆である。
実際 $K$ に対 $30,31$ を両側で追加した六行六列は
$R_0,R_2$ をいずれも初回直線へ制限する完全六空間の座標だからである。
ゆえに $N$ の列空間は座標 $0,1$ への射影が同型な二空間で、座標ベクトル2を含まない。

一方、第二の切断で $K$ を消した残りの全列は、$R_1$ の四座標 $10,11,20,21$ が零である
二次元核を張る。この核の元の $a3$ 成分の全体は $N$ の列空間そのものである
（残りの全階数が2で、既に $N$ が階数2だからである）。
$F_0$ の中心縮約から、$R_1$ には $L_1$ の消滅空間、すなわち行 $2,3$ だけに支持される三関係がある。
それらの三座標制限は式\eqref{eq:pure-branch-core-four}の $20,21$ の二座標だけなので、
非零な一関係がこの核に入る。その $a3$ 成分は座標ベクトル2の倍で、零なら一行3だけに支持されて純関係になる。
従って $N$ の列空間が座標ベクトル2を含み矛盾する。
\end{proof}

\begin{proof}[三中心の排除]
位置0だけは全座標を保ち、他の三位置で三座標制限を行う。
このテンソルを $C(a,b)$ で評価すると全係数が一つの二項一次式で割れる。
実際、二つの内側の対には式\eqref{eq:pure-branch-core-four}を課し、
位置 $0,1$ では $F_0\times L_1$ と $C\to M_1$、位置 $3,0$ では $M_3\times B_0$ と
$L_3\to C$ の全係数を消す。得られる解の基底は、同じ二次式テンソルに $a,b$ をそれぞれ掛けた二本である。
添付計算では四つの三次係数を一座標ずらすことでこれを確認する。零テンソルの場合を含め、
三座標制限が全て消える点 $(a,b)$ を取れる。まず $a\ne0$ とし $u=b/a$ とする。

$U_1,U_2$ を $R_1,R_2$ の両文字の三座標制限の核とする。両核は二次元である。
像次元は式\eqref{eq:pure-branch-core-four}により高々4で、$R_1$ を $L_2$ に、
$R_2$ を $M_2$ に制限した完全六座標のうちそれぞれ四つは三座標だけで読めるからである。
各核から、第四行の最初の三成分を読む写像と、第四列の最初の三成分を読む写像はいずれも単射である。
消えれば第四列だけ、または第四行だけで純テンソルになるためである。
とくに $U_1$ の第四行像は座標ベクトル2を含まない。この向きの $L_2$ への制限が単射だからである。

$w(C(1,u),-,-,-)$ を $v_3$（位置3の文字）の最初の三座標で縮約すると全て $U_1$ に入る。
これらは少なくとも一次元を張る。全て消えれば位置3の支持が高々1となるからである。
その第四行を集めた位置 $2,3$ の三行三列 $H$ は従って非零である。
$H$ は $R_2$ の三座標制限なので式\eqref{eq:pure-branch-core-four}の形であり、
$C(1,u)$ の前点 $L_3(1,u)$ を消すので $d=cu, h=-fu$。
$c=0$ なら非零列は座標ベクトル2の倍に限られ、先の $U_1$ の第四行像に反する。
従って $c\ne0$ で $H$ は階数2である。

同じ $U_1$ の縮約の第四列を集め、位置2を3に固定した位置 $1,3$ の三行三列を $Z$ とする。
核からの二つの読み方が単射なので $Z$ も階数2である。逆に位置1の三座標で縮約すれば全て $U_2$ に入る。
従って $Z$ の行空間は $U_2$ の第四行像の二空間そのものである。
この像は座標ベクトル0を含む。実際、$B_0$ の中心縮約は $R_2$ に $M_3$ を消す列 $0,3$
だけの三関係を与え、三座標制限で $10,20$ の二座標しか持たず、一関係は核に入る。
その第四行の最初の三成分は非零な座標0の倍である。一方 $Z$ の全行は前点
$L_3(1,u)=(1,u,0,0)$ を消すので座標0を含み得ない。これで有限の根は排除された。
$a=0$ の場合は語順を反転し位置 $1,3$ を交換する。この配置は再び同じ三中心配置で、
位置0の順・逆中心が交換される。二項座標を交換すればその根は有限になる。
三座標部分は二直線の線型包そのもので語順の反転によって保たれるから、同じ矛盾を得る。
\end{proof}

\begin{corollary}\label{cor:positive-dimensional-pure-free-decision}
純関係のない標準有限Frobenius候補で、一辺でも点対応の射影が正次元なら、
全次数の標準Koszul左右AS正則性と、正則な候補間の添字保持同型は有限に判定できる。
\end{corollary}
\begin{proof}
二次元以上なら定理\ref{thm:no-moving-plane-pencils}, \ref{thm:isolated-surface-fibers}の被覆判定を使う。
全射影が高々一次元なら、二点以上を持つ射影ファイバーの有無を評価の小行列式で調べる。
一つでもあれば本定理により非正則、なければ系\ref{cor:all-one-dimensional-point-graphs}に帰着する。
\end{proof}
